\section{Asymmetric perturbations}
\label{app:variants}

The main text uses the fully symmetric case, Eq.~\ref{eq:dmatrix}, in which both
groups favour their own and are hostile to the other. Five asymmetric variants are
implemented, and they are the natural place to ask which structural ingredient the
transition actually needs.

\begin{table}[h]
  \centering
  \small
  \begin{tabular}{@{}clll@{}}
    \toprule
    & $D / d$ & \textbf{Reading} & \textbf{Question it answers} \\
    \midrule
    $D_1$ & $\bigl(\begin{smallmatrix} +1 & 0 \\ 0 & 0\end{smallmatrix}\bigr)$
      & one group, in-group tolerance only
      & is in-group favouritism alone enough? \\
    $D_2$ & $\bigl(\begin{smallmatrix} 0 & -1 \\ 0 & 0\end{smallmatrix}\bigr)$
      & one group, out-group intolerance only
      & is out-group hostility alone enough? \\
    $D_3$ & $\bigl(\begin{smallmatrix} +1 & -1 \\ 0 & 0\end{smallmatrix}\bigr)$
      & one group, both
      & does one perturbed group suffice? \\
    $D_4$ & $\bigl(\begin{smallmatrix} +1 & 0 \\ 0 & +1\end{smallmatrix}\bigr)$
      & both groups, in-group tolerance only
      & \\
    $D_5$ & $\bigl(\begin{smallmatrix} 0 & -1 \\ -1 & 0\end{smallmatrix}\bigr)$
      & both groups, out-group intolerance only
      & \\
    $D_6$ & $\bigl(\begin{smallmatrix} +1 & -1 \\ -1 & +1\end{smallmatrix}\bigr)$
      & both groups, both \emph{(main text)}
      & \\
    \bottomrule
  \end{tabular}
\end{table}

\paragraph{Sign convention.} Rows index the receiver's label, columns the
emitter's, and the signs follow \S\ref{sec:spurious}: a positive entry shifts $\hw$
toward agreement, hence toward trust. So $+1$ on the diagonal is in-group
tolerance. \textbf{This is the opposite of the sign printed in the source draft's
Table~I}, whose Eq.~25 and table read literally place the discriminatory phases at
$d < 0$ while its text and all of its figures place them at $d > 0$. One global
sign reconciles the two, and the reading above is the one the sign of $F_\mu$
forces. Fix this before the phase diagram becomes a headline figure --- the $x$
axis of the main result cannot be ambiguous.

\todo{report the four correlation maps for $D_1$--$D_5$. The two rows worth
discussing in text are $D_1$ vs.\ $D_2$ (does tolerance or hostility drive it?) and
$D_3$ vs.\ $D_6$ (does the effect need both groups perturbed, or does one suffice
to organise the whole population?). The second is the more interesting question and
neither the source draft nor the main text here answers it.}

\begin{figure}[h]
  \centering
  \includegraphics[width=0.9\textwidth]{sign_convention_comparison}
  \caption{The two readings of the perturbation's sign, side by side, at otherwise
    identical parameters. Included so that a reader coming from the source draft
    can see immediately which convention each figure uses.}
  \label{fig:signs}
\end{figure}
